\chapter{Main Theorems}

\section{Entropy Non-Amplification}

\begin{theorem}[Capacity accumulation]
Let \((R_1,\ldots,R_T)\) be refinement operators with transcript capacity bounds \(C_1,\ldots,C_T\).  Then
\[
I(X;T_1,\ldots,T_T) \le \sum_{t=1}^{T} C_t .
\]
\end{theorem}

\begin{proof}
Apply the transcript chain rule and bound each conditional term by \(C_t\).
\end{proof}

\section{Depth Lower Bound}

\begin{theorem}[Capacity-forced depth]
If a target has entropy demand \(D=\Theta(n)\) and every admissible refinement step exposes at most \(O(1)\) information, then any complete refinement sequence has depth \(\Omega(n)\).
\end{theorem}

\begin{proof}
A sequence of \(T\) steps exposes at most \(O(T)\) information.  Complete resolution requires \(\Theta(n)\) information, so \(T=\Omega(n)\).
\end{proof}

\section{Operational Rigidity}

\begin{theorem}[Capacity obstruction]
If the required uncertainty reduction exceeds the admissible transcript capacity of every allowed refinement sequence, then stabilization is impossible inside that admissible class.
\end{theorem}

\begin{proof}
Stabilization would require a resolving transcript.  The capacity bound prevents such a transcript from carrying the required information.
\end{proof}

\section{Boundary}

These are framework theorems.  Domain-specific conclusions require verified reductions into the hypotheses of the stated theorems.
